\section{Scope, conventions, and the main theorem}
\label{sec:scope}

The formulas in the source pages are normally stated over ordinary real
numbers and use total-looking notation for quotients, inverse functions and
square roots.  Constructively these are not all statements of the same kind.
We therefore separate a finite algebraic kernel from the analytic presentation
layer.

\begin{definition}[Proof levels]
A formula in this report is one of the following.
\begin{enumerate}[label=\textup{(\roman*)}]
\item A \emph{polynomial identity}: after the displayed abbreviations are
expanded, it is an equality in a commutative ring.  Such a theorem needs no
order, excluded middle or real completeness.
\item A \emph{guarded rational identity}: all denominators carry explicit
apartness witnesses.  Its cross-multiplied form is polynomial.
\item A \emph{positive-root identity}: the radicand is proved nonnegative and
an interval or branch certificate selects the nonnegative square root.
\item A \emph{limit or differential identity}: every derivative, integral or
limit is represented by a fixed primitive-recursive rational-name program and
an explicit primitive-recursive modulus.
\end{enumerate}
\end{definition}

\begin{convention}[Proper oriented triangle]
Unless stated otherwise, \(A,B,C\in S^2\) are three pairwise non-antipodal
points joined by the shorter great-circle arcs.  Their side lengths satisfy
\[
 0<a,b,c<\pi,\qquad a+b+c<2\pi,
\]
and their interior angles satisfy \(0<A,B,C<\pi\).  We choose the orientation
so that the scalar triple product introduced in Section~\ref{sec:vector-model}
is positive.  In the arithmetic interpretation, each strict inequality is
positive data: a rational separation interval, not merely the negation of an
equality.
\end{convention}

The symbol \(q\) denotes a quarter turn and \(\pi=2q\).  The constructive
trigonometric packet cited in the bibliography constructs \(q\) as the unique
zero of cosine in \([1,2]\) by bounded sign search and proves
\(\sin q=1\), \(\cos q=0\).  Consequently no theorem below presupposes a
classical definition of \(\pi\).

\begin{theorem}[Constructive spherical-trigonometry atlas]
\label{thm:main-atlas}
For every fixed tuple of primitive-recursive regular dyadic names encoding a
proper oriented spherical triangle, together with fixed unit-vector,
orientation and apartness certificates, \(\ipra\) proves:
\begin{enumerate}[label=\textup{(\alph*)}]
\item the three spherical side cosine laws, the three angle cosine laws and the
spherical sine law;
\item the polar-triangle correspondences and involution;
\item all cyclic instances of the half-angle, half-side, half-tangent,
haversine, Delambre--Gauss, Napier, tangent, cotangent, five-part and Cagnoli
identities stated in this report;
\item the ten right-triangle and ten quadrantal-triangle rules, with explicit
guards on quotient forms;
\item the differential variation formulas for every fixed primitive-recursive
differentiable one-parameter family of proper triangles;
\item Girard's excess formula and the polygonal excess formula for every fixed
finite geodesic polygon with a certified triangulation;
\item the tetrahedral face-area cosine law and the tetrahedral sine-product
identity for every fixed nondegenerate tetrahedron with oriented face data.
\end{enumerate}
Every output proof can be generated from the finite input programs and their
certificates by a primitive-recursive metatheoretic proof transformer.
\end{theorem}

The quantification in Theorem~\ref{thm:main-atlas} is metatheoretic.  First-order
\(\ipra\) has numerical variables but neither a real-number sort nor function
variables.  Thus the theorem is a schema over fixed primitive-recursive terms;
it is not an assertion that \(\ipra\) contains a universal evaluator for all
primitive-recursive programs.  Section~\ref{sec:ipra} makes the translation
precise, and Appendix~\ref{app:ipra-target} displays a representative
object-language target.

\subsection{What is genuinely new work in the reduction}

The analytic sine/cosine layer can be reused from ProveIt's constructive
trigonometric-identities packet.  The new spherical contribution is the
following small kernel.
\begin{enumerate}
\item Encode a triangle by three unit vectors and retain the scalar triple
product as an orientation witness.
\item Derive all primary laws from two polynomial identities: Lagrange's
cross-product identity and the Gram determinant.
\item Use normalized face normals to construct the polar triangle without a
choice of coordinates.
\item Treat every familiar quotient formula as a cross-multiplied identity
plus explicit apartness data.
\item For tetrahedra, replace edge vectors by the four oriented face-area
vectors; their sum is zero, so the advertised cosine law is an ordinary norm
expansion.
\end{enumerate}
This organization makes the mathematical and formalization dependencies
transparent: the bulk of spherical trigonometry is finite algebra, not real
analysis.
